\chapter{Model Family Overview}
\label{chap:models}

We now will walk though the intuition behind some model families.
Future Chapters dive into some specific model families in a \textit{lot} of detail - these will be skipped here.
Generally, model types in this chapter are outdated, but are kept alive by the \textsc{SciKit-Learn} fans and Coursera\textsuperscript{TM}. 
But first, a word about \textit{why} we need to impose assumptions on our models. 


\section{No Free Lunch}
\label{chap:nfl}

\subsection{Finite Observations Do Not Generalise}
\label{sec:nfl-why-finite-observations-are-not-enough}
The observations $(1,2)$, $(2,4)$ and $(3,6)$ fit both
\[
    f_1(x)=2x,
    \qquad
    f_2(x)=2x+(x-1)(x-2)(x-3),
\]
yet the functions disagree at unseen inputs.
Intravitally, we should choose the simpler, linear function.
However, this is an assumption we are imposing on our model, and we should remember it.
\figpdf[.79\textwidth]{nfl_two_functions}{Two functions can fit the same observations perfectly and disagree elsewhere; prediction requires a preference among compatible hypotheses.}

\subsection{Occam's Zazor}
\label{sec:nfl-occam-s-razor-is-also-a-bias}
``Prefer simple explanations'' is useful only after choosing what simplicity means.
A periodic physical law can be compact in a Fourier representation and awkward in a polynomial one.

\subsection{NFL is Compatible with our Toolbox}
NFL tells us that \textit{to get better task-specific accuracy we need to add assumptions}.
We add assumptions with almost all of our design choices:
\begin{itemize}
    \item \textbf{Architecture:} Data modality and task assumptions.
    \item \textbf{Statistical Guarantees:} sampling process, a restricted hypothesis family, a loss, or a confidence statement.
    A useful intuition:
    \[
        \underbrace{\text{generalisation gap}}_{\text{population error} - \text{training error}}
        \sim
        \sqrt{\frac{\underbrace{\text{model-class complexity}}_{\text{more hypotheses cost more evidence}}+\underbrace{\text{confidence penalty}}_{\text{stronger guarantee costs evidence}}}{N_{\text{examples}}}}.
    \]
    The exact theorem changes with the assumptions and complexity measure.
    The important point is that guarantees are conditional on structure; they do not contradict an assumption-free averaging theorem.
    \item \textbf{Optimisers:} NFL results average search algorithms over broad objective-function families under their own assumptions.
    Practical optimisers exploit gradients, curvature, smoothness and architecture-specific structure, so they are not being compared under an assumption-free real-world distribution.
    The optimiser evolution in \cref{sec:optimisation-how-optimiser-ideas-evolved-a-branching-genealogy} is therefore a sequence of useful structural bets, not a contradiction of NFL.
    \item \textbf{Domain:} Symmetries, topology, conservation laws, known dynamics, simulator structure or measurement constraints.
\end{itemize}



\section{Linear}
\label{sec:models-linear-models-a-transparent-baseline}
\paragraph{Linear regression} predicts with a feature-wise score:
\[
    \hat y=\underbrace{\theta_{\text{weights}}^\top x}_{\text{weighted features}}+\underbrace{\theta_{\text{bias}}}_{\text{offset}}.
\]
``Linear'' refers to the chosen features.
Adding nonlinear engineered features can create nonlinear relationships in raw-input space while the model remains linear in its parameters.
\paragraph{Logistic regression} passes the same score through a sigmoid to model a binary probability.
\[
    \hat y = \frac{1}{1 + e^{-(\theta_{\text{weights}}^\top x + \theta_{\text{bias}})}}
\]



\section{Support Vector Machines and Kernels}
\label{sec:models-support-vector-machines-and-kernels}
Support Vector Machine (SVM) draws a decision line between classes in the feature space.
We aim to maximise the margin between the \textit{Support Vectors}, which are the closest datapoints either side of the line, and the decision line. 
This decision line is a hyperplane in an N-dimensional space, but just imagine it as drawing lines.


\paragraph{Linear vs. Non-linear SVM}
Linear SVM draws a straight decision boundary in the original feature space. Non-linear SVM uses kernels to construct curved decision boundaries, allowing classes that cannot be separated by a straight line to be distinguished.


\paragraph{Hard vs. Soft}
Hard SVM requires decision lines to exactly differentiate between classes, while soft SVM allows for overlap by penalising overlap weakly with a regularisation parameter.
A linear soft-margin SVM trades margin size against hinge-loss violations:
\[
    \min_{\theta_{\text{weights}},\theta_{\text{bias}}}
    \underbrace{\frac12\|\theta_{\text{weights}}\|_2^2}_{\text{prefer a wider margin / smaller weights}}
    +
    \underbrace{C\sum_{i=1}^{N_{\text{examples}}}
    \max\!\left(0,1-y_i(\theta_{\text{weights}}^\top x_i+\theta_{\text{bias}})\right)}_{\text{penalise margin violations}}.
\]
The first term is an L2-style parameter preference, directly linking the SVM objective to \cref{sec:stability-regularisation-l1-and-l2}.
Larger $C$ makes training violations more expensive relative to margin width.

\paragraph{Support Vector Regression (SVR)}
SVR applies the same principle to continuous targets rather than classes.
Instead of separating classes, we fit a function with an $\epsilon$-wide tolerance around it. 
Predictions within this tolerance incur no loss, while errors outside it are penalised.


\paragraph{Kernels}
When data cannot be separated by a flat line (e.g., one group of points forms a circle inside another group), SVM uses kernel functions (such as Polynomial or Radial Basis Function / RBF kernels). 
When temporarily working in a different feature space, the possibility of drawing a differentiating line is higher. 
The \textit{kernel trick} mathematically projects the data into a higher-dimensional space where a linear hyperplane can separate them.

A kernel replaces explicit feature construction with an inner product in a intermediate feature space:
\[
    \underbrace{k(x,x')}_{\text{similarity used by the model}}
    =
    \left\langle
    \underbrace{\phi(x)}_{\text{implicit features}},
    \phi(x')
    \right\rangle.
\]
A linear separator in that feature space can therefore become nonlinear in the original coordinates.
\figpdf[.9\textwidth]{kernel_feature_map}{A nonlinear boundary in raw coordinates can become simple after a feature transformation; kernels can compute some such similarities without explicitly materialising the feature map.}



\section{Neural Networks}
\label{sec:models-neural-networks-and-nonlinear-activations}
You will hear a lot about Neural networks, and their descendants (RNNs, CNNs, etc.) and I will now introduce their smallest componant. 
Their \textit{LEGO bricks}, from which they are all built from.

\subsection{Neuron}
A neuron is the fundamental computational unit of a neural network.
It receives input values, multiplies by a learned weight, adds a bias, and passes the result through an \textit{activation function}.
\[
    \underbrace{h}_{\text{neuron output}}
    =
    \underbrace{\phi}_{\text{activation}}
    \left(
        \underbrace{\sum_{i=1}^{d} \theta_i x_i}_{\text{weighted inputs}}
        +
        \underbrace{\theta_{\text{bias}}}_{\text{offset}}
    \right).
\]
The weights determine how strongly each input influences the output, while the activation function determines how the combined signal is transformed.

\paragraph{Activation Functions}
The activation functions are where Neural Networks derive their non-linearity.
Sigmoid/tanh saturate at large magnitude, ReLU keeps a constant positive-side derivative but can create inactive units, GELU gives a smooth input-dependent gate-like response.
Activation choice interacts with initialisation and gradient scale in \cref{sec:stability-xavier-glorot-and-he-kaiming-initialisation}.
\figpdf[.81\textwidth]{activations}{Common nonlinearities differ mainly in saturation, zeroing behaviour and local derivative scale.}

\begin{itemize}
    \item \textbf{Linear (Identity):} Returns the input unchanged. Common in regression output layers. $f(x)=x$.

    \item \textbf{Binary Step:} Outputs $0$ or $1$ depending on whether the input crosses a threshold. Used in classical perceptrons, but unsuitable for ordinary gradient-based training. $f(x)=\begin{cases}1 & x\geq0\\0 & x<0\end{cases}$.

    \item \textbf{Sigmoid:} Squashes values into $(0,1)$, useful for binary classification probabilities and gating mechanisms. Can suffer from vanishing gradients. $f(x)=\frac{1}{1+e^{-x}}$.

    \item \textbf{Tanh:} Squashes values into $(-1,1)$, providing zero-centred outputs. Historically common in RNNs and LSTMs. $f(x)=\tanh(x)$.

    \item \textbf{ReLU (Rectified Linear Unit):} Sets negative values to zero and leaves positive values unchanged. Computationally cheap and widely used in neural networks, but can produce permanently inactive neurons. $f(x)=\max(0,x)$.

    \item \textbf{Leaky ReLU:} Like ReLU, but allows a small negative slope to reduce the risk of inactive neurons. $f(x)=\max(\alpha x,x)$, where $0<\alpha<1$.

    \item \textbf{PReLU (Parametric ReLU):} Like Leaky ReLU, but learns the negative slope $\alpha$ during training. $f(x)=\max(\alpha x,x)$.

    \item \textbf{ELU (Exponential Linear Unit):} Behaves like ReLU for positive inputs but smoothly saturates towards $-\alpha$ for negative inputs. $f(x)=\begin{cases}x & x>0\\\alpha(e^x-1) & x\leq0\end{cases}$.

    \item \textbf{SELU (Scaled ELU):} A scaled ELU designed to encourage self-normalising behaviour under specific architectural and initialisation conditions. $f(x)=\lambda\begin{cases}x & x>0\\\alpha(e^x-1) & x\leq0\end{cases}$.

    \item \textbf{GELU (Gaussian Error Linear Unit):} Smoothly gates inputs according to their magnitude, rather than applying ReLU's hard threshold. Common in Transformers such as BERT and GPT-2. $f(x)=x\Phi(x)$, where $\Phi$ is the standard normal CDF.

    \item \textbf{SiLU (Sigmoid Linear Unit / Swish):} Smooth activation that multiplies the input by its sigmoid. Used in architectures such as EfficientNet and Llama's SwiGLU feed-forward layers. $f(x)=x\sigma(x)$.

    \item \textbf{Swish:} Generalised SiLU with an adjustable parameter $\beta$ controlling the sigmoid gate. $f(x)=x\sigma(\beta x)$.

    \item \textbf{Mish:} Smooth, non-monotonic activation combining the input with a softplus transformation. $f(x)=x\tanh(\ln(1+e^x))$.

    \item \textbf{Softplus:} Smooth approximation to ReLU that never produces exactly zero for finite inputs. $f(x)=\ln(1+e^x)$.

    \item \textbf{Softsign:} Smoothly compresses values into $(-1,1)$, with slower saturation than tanh. $f(x)=\frac{x}{1+|x|}$.

    \item \textbf{Softmax:} Converts a vector of logits into a probability distribution whose elements sum to $1$. Common in multiclass classification and attention. $f(x)_i=\frac{e^{x_i}}{\sum_j e^{x_j}}$.

    \item \textbf{GLU (Gated Linear Unit):} Multiplies one learned projection by the sigmoid of another, allowing the network to gate information. $f(x)= (W_a x+b_a)\odot\sigma(W_b x+b_b)$.

    \item \textbf{SwiGLU:} Gated activation combining SiLU with another learned projection. Used in modern Transformer feed-forward networks, including Llama. $f(x)=\operatorname{SiLU}(W_a x+b_a)\odot(W_b x+b_b)$.
\end{itemize}


\subsection{Layers}
We combine many neurons into layers and stack them in sequence.
This gives interesting dimensions of variability, because we can:
\begin{itemize}
    \item \textbf{Add Non-Neuron functions} like Pooling, Attention, Dropout, etc.
    \item \textbf{Loop Layers} by passing outputs back as inputs,
    \item \textbf{Ensure Statistical Properties} like normalisation, 
    \item \textbf{Vary Width \& Depth} Vary the number of neurons per layer and the number of layers in the network 
    \item \textbf{Connect Layers Differently} Skip layers, branch into parallel paths, or combine multiple representations.
    \item \textbf{Share Weights} Reuse the same learned parameters across different positions or time steps.
\end{itemize}

This gives a lot of dimensions researchers can explore, and \textit{many} permutations to investigate.
Heuristically, the study of neural networks is either adding to this list through discovery or experimentally trial-and-error-ing a new permutation.

\subsubsection{Common Layer Types}
\begin{itemize}
    \item \textbf{Dense / Fully Connected:} Connects every input feature to every output neuron. Common in Multi-Layer Perceptrons (MLPs) and classification heads.
    \item \textbf{Convolutional:} Applies small, shared filters across local regions of an input. Common in CNNs for images, audio, and other structured signals.
    \item \textbf{Pooling:} Reduces the spatial or sequence dimensions by aggregating values, commonly using maximum or average pooling.
    \item \textbf{Recurrent:} Reuses a state-update operation across sequence steps, allowing previous information to influence later outputs. Used in RNNs, LSTMs, and GRUs.
    \item \textbf{Embedding:} Maps discrete identifiers, such as token IDs, to learned continuous vectors.
    \item \textbf{Attention:} Computes weighted combinations of representations based on their relevance to one another. Central to Transformers.
    \item \textbf{Normalisation:} Rescales or normalises intermediate activations to improve training behaviour. Examples include Batch Normalisation, Layer Normalisation, and RMSNorm.
    \item \textbf{Dropout:} Randomly masks activations during training to reduce overfitting. Typically disabled during inference.
    \item \textbf{Residual / Skip Connection:} Adds an earlier representation to a later one, allowing information and gradients to bypass intermediate transformations.
    \item \textbf{Gated:} Controls information flow using learned multiplicative gates. Common in LSTMs, GRUs, and modern Transformer feed-forward blocks.
    \item \textbf{Upsampling:} Increases the spatial resolution of a representation. Common in image generation, segmentation, and decoder architectures.
    \item \textbf{Output / Projection:} Maps an internal representation into the required output dimension, such as class logits or vocabulary scores.
\end{itemize}

\section{Model assumptions and inductive biases}
\label{sec:models-model-assumptions-and-inductive-biases}
\Cref{sec:nfl-inductive-bias-the-assumptions-that-make-learning-possible} 
\small
\begin{longtable}{@{}L{.22\linewidth}L{.44\linewidth}L{.25\linewidth}@{}}
\caption{Model families, their structural preference, and a representative useful regime.}\label{tab:model_biases}\\
\toprule
Model & Structural preference / assumption & Potential useful setting \\
\midrule
\endfirsthead
\toprule
Model & Structural preference / assumption & Potential useful setting \\
\midrule
\endhead
Linear / logistic regression & Chosen features combine additively in an affine score & Interpretable baseline effects \\
KNN & Nearby observations should have similar targets & Local smooth structure \\
Kernel SVM & A chosen similarity creates a useful separating feature space & Nonlinear boundaries, moderate data \\
Decision tree & Hierarchical thresholds capture useful interactions & Heterogeneous tabular relations \\
Random forest & Averaging decorrelated trees reduces instability & Nonlinear tabular baselines \\
Gradient boosting & Sequential corrections improve a chosen objective & Structured/tabular data \\
PCA & High-variance linear directions are useful & Compression/exploration \\
Gaussian process & Kernel covariance describes function similarity & Small-data regression with uncertainty \\
MLP & Learned nonlinear compositions capture the relation & General vector inputs \\
CNN & Local patterns repeat across positions & Images/gridded signals \\
RNN / GRU / LSTM & Recurrent state summarises useful past context & Sequences/streaming \\
Transformer & Content-dependent token interactions matter & Context-rich sequence/multimodal tasks \\
GNN & Graph connectivity defines relevant neighbours & Molecules/relational systems \\
SSM / Mamba & Long history can be compressed into evolving state & Long or streaming sequences \\
Autoencoder / VAE & Reconstruction preserves useful latent structure & Compression/generative latents \\
GAN / diffusion / flow & Different generative mechanisms match complex data distributions & Sample/density modelling \\
Mixture of experts & Different inputs benefit from specialist subnetworks & Large conditional-capacity systems \\
Siamese / contrastive & Related examples should be close in learned space & Retrieval/few-shot tasks \\
\bottomrule
\end{longtable}
\normalsize
